\documentclass[11pt, a4paper]{article}

\usepackage{fontspec}
\setmainfont{Latin Modern Roman}
\setmonofont{DejaVu Sans Mono}[Scale=0.85]
\usepackage[margin=2.5cm]{geometry}
\usepackage{amsmath, amssymb, amsthm}
\usepackage{xurl}
\usepackage{hyperref}
\usepackage[dvipsnames]{xcolor}
\usepackage{enumitem}
\usepackage{fancyvrb}
\usepackage{caption}
\usepackage{draftwatermark}
\SetWatermarkText{PREPRINT}
\SetWatermarkScale{1.1}
\SetWatermarkLightness{0.92}
\captionsetup[figure]{name=Listing}
\emergencystretch=5em
\newcommand{\brk}{\discretionary{}{}{}}
\newcommand{\lean}[1]{\ifmmode\text{\ttfamily\detokenize{#1}}\else{\ttfamily\detokenize{#1}}\fi}

\definecolor{indigo}{rgb}{0.29, 0.0, 0.51}
\definecolor{teal}{rgb}{0.0, 0.5, 0.5}
\hypersetup{colorlinks=true, linkcolor=teal, citecolor=indigo, urlcolor=blue,
  pdftitle={Fricke-self-dual eta-quotients (preprint)}, pdfauthor={Xavier Callens}}

\newtheorem{theorem}{Theorem}[section]
\newtheorem{lemma}[theorem]{Lemma}
\newtheorem{definition}[theorem]{Definition}
\newtheorem{corollary}[theorem]{Corollary}
\theoremstyle{remark}
\newtheorem{remark}[theorem]{Remark}

\title{\textbf{Fricke-Self-Dual Eta-Quotients Are Eigenforms}\\[0.3em]
  \large\textbf{A Corollary in Lean 4, and a Note on Its Motivation \\ in Persson--Volpato's CHL S-Duality}\\[0.8em]
  \large\textcolor{BrickRed}{PREPRINT --- not peer reviewed}}
\author{
  Xavier Callens\thanks{Corresponding author: \texttt{callensxavier@gmail.com}} \\
  \small SocrateAI Research Initiative
}
\date{\today\\[0.6em]
  \small Archived at Zenodo: \texttt{ZENODO\_DOI\_PLACEHOLDER}}

\begin{document}
\maketitle

\begin{center}
\fbox{\begin{minipage}{0.92\textwidth}
\small \textbf{Status.} This is a \emph{preprint}. It has \textbf{not been peer reviewed}. The
Lean development it describes compiles and is axiom-audited --- those claims are machine-checked
and independently reproducible --- but the \emph{exposition} has had no external referee. The
mathematical content is a single closed corollary (Theorem~\ref{thm:main}) of an already-proved
transformation law; the physics motivation cited in \S\ref{sec:motivation} is reported honestly as
\emph{motivation}, not as something this paper formalizes or proves. Read the scope before citing
the result.
\end{minipage}}
\end{center}

\begin{abstract}
We formalize, in Lean~4 on top of Mathlib, a corollary of an already-proved Fricke transformation
law for eta-quotients: if the exponent vector $r$ of $f(z) = \prod_{\delta\mid N}\eta(\delta z)^{r_\delta}$
satisfies $r_\delta = r_{N/\delta}$ for every $\delta\mid N$ (\emph{Fricke-self-duality}), then $f$
is a genuine eigenform of the Fricke transform $z\mapsto -1/(Nz)$, with a closed-form eigenvalue
$\lambda = i^{-k}N^{k/2}$ depending only on the level $N$ and weight $k$, not on $r$. We also prove
exactly when $\lambda=\pm1$. The only new mathematical content is a single product identity,
$\prod_{\delta\mid N}\delta^{r_\delta}=N^k$ on self-dual $r$; every other declaration is a direct
substitution into the pre-existing transformation law. We report, as motivation and at the
literature tier of this programme's epistemic ladder --- not as a formalized claim --- that this
self-duality condition is exactly Persson and Volpato's ``balanced Frame shape,'' which they
identify as characterizing CHL string compactifications self-dual under a Fricke-type S-duality on
the heterotic axio-dilaton modulus. We state precisely where their normalization and ours
diverge, so the two are not conflated.
\end{abstract}

\noindent\textbf{MSC 2020:} 11F11, 11F20, 68V20.\quad
\textbf{Keywords:} eta-quotients, Fricke involution, eigenforms, Lean 4, Mathlib, formal
verification, CHL models.

\section{Introduction}\label{sec:introduction}

A companion development in this programme, \emph{A Mathlib-Idiomatic Formalization of the Fricke
Involution}~\cite{Fricke2026}, formalizes the matrix-level Fricke involution $W_N$ on
$\Gamma_0(N)$, and a second, \emph{Eta-Quotients in Lean 4}~\cite{EtaQuotients2026}, formalizes the
transformation law of eta-quotients under $W_N$: for $f(z)=\prod_{\delta\mid N}\eta(\delta
z)^{r_\delta}$ of weight $k$ (i.e.\ $\sum_{\delta\mid N} r_\delta = 2k$),
\begin{equation}\label{eq:etaQuotient_fricke}
  f\!\left(\frac{-1}{Nz}\right) \;=\; i^{-k}N^k\Bigl(\sqrt{\textstyle\prod_{\delta\mid N}\delta^{r_\delta}}\Bigr)^{-1} z^k \, f^\ast(z),
\end{equation}
where $f^\ast$ is the \emph{dual} quotient, exponent vector $\delta\mapsto r_{N/\delta}$. This is
\lean{etaQuotient_fricke} (tag \texttt{F3.2-A7}), already proved, sorry-free, in
\texttt{EtaQuotientModularity.lean}. In general $f^\ast\neq f$: \lean{fricke_dual_pin} in that
module exhibits an $r$ whose dual is a genuinely different function.

This note asks the natural next question: what happens when $f^\ast=f$? Equation
\eqref{eq:etaQuotient_fricke} then collapses to a genuine eigenform relation, with an eigenvalue
depending only on $N$ and $k$. That collapse is short --- a substitution and a simplification, not
a fresh derivation --- and its only nontrivial ingredient is one product identity
(Lemma~\ref{lem:prod-collapse}). We formalize it, together with a closed characterization of when
the eigenvalue is exactly $\pm1$, and we record, honestly and at arm's length, why the question was
worth asking: the self-duality condition on $r$ is exactly the combinatorial shape of a
condition in the physics literature that this paper does not attempt to derive.

\section{The self-duality condition}

\begin{definition}[Fricke-self-dual exponent vector]\label{def:selfdual}
\lean{IsFrickeSelfDual}. For $N>0$ and an exponent vector $r:\mathbb{N}\to\mathbb{Z}$,
\[
  \mathrm{IsFrickeSelfDual}(N,r) \;:\Longleftrightarrow\; \forall\, \delta\mid N,\ r_\delta = r_{N/\delta}.
\]
\end{definition}

We call this condition \emph{self-dual} (equivalently \emph{Fricke-symmetric}), never
``balanced.'' \texttt{EtaQuotientModularity.lean} already uses ``balanced'' for an unrelated
concept --- \lean{exists_balanced_add}, tag \texttt{F3.2-B1}, the \lean{Int.bmod}
balanced-remainder-division step used in the $\Gamma_0(N)$ generator descent --- and the physics
literature discussed in \S\ref{sec:motivation} uses ``balanced'' for a third, related but
distinct, sense. Reusing the word here would collide with both. \lean{IsFrickeSelfDual} is
\lean{Decidable} (\lean{decidableIsFrickeSelfDual}), so membership can be checked by
\lean{decide} on concrete instances, which every pin in the development below does.

\begin{lemma}[The product collapses to a level power]\label{lem:prod-collapse}
\lean{prod_zpow_selfDual}. For $N>0$, $r$ self-dual on $N$, and $\sum_{\delta\mid N} r_\delta=2k$,
\[
  \prod_{\delta\mid N} \delta^{r_\delta} \;=\; N^k.
\]
\end{lemma}

This is the only place the self-duality hypothesis does real work beyond a substitution: it forces
the divisor-weighted product of the exponents to collapse to a single power of the level, which is
what removes $r$ from the eigenvalue entirely. It was checked, independently of the Lean proof, in
exact rational arithmetic over four ranges (up to $1\le N\le 60$, exponents in $[-4,4]$,
$671504$ self-dual instances), zero mismatches, and against $9597$ non-self-dual even-weight
vectors of which $98.8\%$ violate the identity --- confirming self-duality is load-bearing for the
collapse, not an artifact of the weight condition alone.

\section{The eigenform relation}

\begin{definition}[Fricke eigenvalue]\label{def:eigenvalue}
\lean{frickeEigenvalue}. For $N>0$ and $k\in\mathbb{Z}$,
\[
  \lambda(N,k) \;:=\; i^{-k}\sqrt{N^{k}}.
\]
\end{definition}

\begin{theorem}[Self-dual eta-quotients are Fricke eigenforms]\label{thm:main}
\lean{etaQuotient_fricke_selfDual}. For $N>0$, $r$ self-dual on $N$, $\sum_{\delta\mid N}
r_\delta=2k$, and $z$ in the open upper half-plane,
\[
  f\!\left(\frac{-1}{Nz}\right) \;=\; \lambda(N,k)\, z^k\, f(z), \qquad f(z)=\prod_{\delta\mid N}\eta(\delta z)^{r_\delta}.
\]
\end{theorem}

The proof is two rewrites: substitute \eqref{eq:etaQuotient_fricke}'s dual quotient
$f^\ast$ by $f$ itself (Definition~\ref{def:selfdual} makes this a congruence over the divisor
product, \lean{etaQuotient_congr_divisors}), then collapse the constant using
Lemma~\ref{lem:prod-collapse} (\lean{fricke_const_selfDual}). Nothing else enters. An
intermediate, unnormalized form retaining $\sqrt{\prod\delta^{r_\delta}}$ rather than $\lambda$ is
\lean{etaQuotient_fricke_selfDual_raw}; a normalized form isolating the root of unity
$i^{-k}$ on one side is \lean{etaQuotient_fricke_selfDual_normalized}.

\begin{remark}[What this is not]\label{rem:not-modularity}
Theorem~\ref{thm:main} makes no reference to $\Gamma_0(N)$, to cusp conditions, or to the bundled
type $M_k(\Gamma_0(N))$: it is a statement about one function value at one point,
$-1/(Nz)$, transported from the already-proved general transformation law. Calling
$\lambda(N,k)$ ``the eigenvalue of the Fricke operator $W_N$ on $M_k(\Gamma_0(N))$'' would overstate
what is proved here; the correct reading is that $f$ is an eigenvector of the single linear map
$g\mapsto g(-1/(Nz))$ restricted to functions of this shape.
\end{remark}

\subsection{When is the eigenvalue $\pm1$?}

The eigenvalue's dependence on $N$ and $k$ alone makes an exact characterization of
$\lambda=\pm1$ reachable without further hypotheses on $r$.

\begin{theorem}[Modulus, and the sign criteria]\label{thm:pm1}
For $N>0$ and $k\in\mathbb{Z}$:
\begin{align}
\|\lambda(N,k)\| &= \sqrt{N^{k}} \tag{\lean{frickeEigenvalue_norm}}\\
\lambda(N,k)=1 &\iff 4\mid k \ \wedge\ (N=1 \vee k=0) \tag{\lean{frickeEigenvalue_eq_one_iff}}\\
\lambda(N,k)=-1 &\iff k\bmod 4 = 2 \ \wedge\ N=1 \tag{\lean{frickeEigenvalue_eq_neg_one_iff}}
\end{align}
\end{theorem}

\begin{corollary}[$\pm1$ in terms of the weight sum]\label{cor:pm1-weight}
\lean{frickeEigenvalue_pm_one_iff}. For $N>0$, $r$ self-dual, $\sum_{\delta\mid N} r_\delta=2k$,
\[
  \lambda(N,k)\in\{1,-1\} \iff 4 \Bigm| \textstyle\sum_{\delta\mid N} r_\delta \ \wedge\ \bigl(N=1 \vee \sum_{\delta\mid N} r_\delta = 0\bigr).
\]
\end{corollary}

Theorem~\ref{thm:pm1} rests on two independent arithmetic facts, each proved and pinned on its own
before being combined: the periodicity of $i^{k}\bmod 4$ (\lean{I_zpow_emod}, genuinely new ---
Mathlib's own \lean{Complex.I_pow_eq_pow_mod} is stated only for $\mathbb{N}$-exponents, not
$\mathbb{Z}$), and $(N:\mathbb{R})^k=1 \iff (N=1\vee k=0)$ for $N>0$
(\lean{natCast_zpow_eq_one_iff}). Every pinned instance in both theorems is checked twice, by the
general lemma and by an independent hand computation sharing no lemma with it; the informative
instances are the ones where a naive guess would be wrong --- e.g.\ $(N,k)=(6,2)$, where $4\mid k$
fails but $k\bmod 4=2$ holds and $\lambda\neq -1$ only because $N\neq1$, the case that most directly
tests whether the level condition was dropped by accident.

\section{Verification status}\label{sec:verification}

The development is 333 declarations in 6437 lines, in one new module,
\texttt{EtaQuotient\brk FrickeSelfDual.lean}, imported into the library's default build target. It is
entirely \lean{sorry}-free: 24 statement-graph nodes (two definitions, sixteen general lemmas and
characterizations, five worked eigenvalue pins cross-checked against a $q$-product evaluation of
$\eta$, and one statement-shape guard), all \texttt{proved}, none blocked, none open. 299 of the
333 declarations carry a build-failing \lean{\#guard\_msgs in \#print axioms} guard reporting
exactly \lean{[propext, Classical.choice, Quot.sound]}; the remainder are the pure data
definitions (the pinned exponent vectors) and auxiliary lemmas that a guard would not
strengthen. \lean{lake build SocrateAI} --- the whole library's shared default target, now also
carrying this module's guards --- completes in 3769 jobs with no errors.

Every headline theorem was checked, independently of its Lean proof, against a direct
$q$-product evaluation of the Dedekind eta function (a 600-term truncation, three points per
instance, relative error $\le 1.8\times10^{-14}$), and Theorem~\ref{thm:main}'s $N=1$,
$r\equiv24$, $k=12$ instance reproduces \lean{eta_S_via_fricke}, itself independently re-derived
from Mathlib's own \lean{discriminant_S_invariant}. Kernel proof terms were inspected directly
(\lean{\#print}) to confirm the headline theorems are genuine substitutions into
\lean{etaQuotient_fricke} rather than restatements that merely typecheck: each reduces to at most
two \lean{rw} steps plus \lean{ring}.

The statement-dependency graph (\lean{dag/theorems.jsonl}) validates: 104 nodes, 95 proved,
acyclic. No axiom beyond Lean's three standard axioms and the library's six pre-existing,
disclosed \lean{physics\_postulate\_$i$} placeholders (unused by anything proved here) is declared
anywhere in the library.

\section{Motivation: Persson--Volpato's Fricke S-duality in CHL models}\label{sec:motivation}\label{sec:selfdual-pointer}

\textbf{Nothing in this section is formalized. Nothing in \S\S\ref{def:selfdual}--\ref{sec:verification} names or
depends on any object introduced here.} We report the motivation for asking the question in the
first place, at the literature (L) tier of this programme's epistemic ladder, because the
disclosure norm this programme follows~\cite{Tao2026AI,Leiden2026} requires attributing sources
plainly rather than letting a formal artifact imply more provenance than it has.

Persson and Volpato~\cite{PerssonVolpato2015} study CHL compactifications --- orbifolds of the
heterotic string (equivalently, by string duality, of type~IIA on $K3\times T^2$) by an order-$N$
symmetry $g$ --- and attach to $g$ a \emph{generalized Frame shape}
$g\leftrightarrow\prod_{a\mid N}a^{m(a)}$, an exponent vector on the divisors of $N$, constrained
by $\sum_{a\mid N}m(a)\cdot a=24$ (their eq.~1.4, recording a $24$-dimensional representation).
They call the Frame shape \emph{balanced} (their eq.~1.6) when $m(N/a)=m(a)$ for all $a\mid N$ ---
exactly Definition~\ref{def:selfdual}, read on their exponent vector --- and their headline result
(their \S1.1, \S3) is that a CHL model is self-dual under a Fricke S-duality $S\mapsto -1/(NS)$ on
the heterotic axio-dilaton modulus $S$ precisely when its Frame shape is balanced. That
identification is established by physical arguments this Lean library does not
contain and this note does not attempt: charge-lattice $N$-modularity, Conway-group
representation theory on the $24$-dimensional representation, and Witten-index/BPS-state counting.

\begin{remark}[Two normalizations that must not be merged]\label{rem:normalization}
Persson--Volpato's weight constraint $\sum_{a\mid N}m(a)\cdot a=24$ is \emph{not} the eta-quotient
weight normalization $\sum_{\delta\mid N}r_\delta=2k$ used throughout this paper. The two agree
only at $N=1$. Their canonical balanced shape at $N=2$, $1^8 2^8$, satisfies $\sum a\cdot m(a) =
8+16=24$ but $\sum m(a)=16$, i.e.\ weight $8$, not $12$ (checked this session over
$1^k N^k$, $k=24/(1+N)$, for $N\in\{1,2,3,5,7,11,23\}$). A reader translating between the two
literatures must carry the exponent vector's actual weight, not assume weight $12$ throughout.
\end{remark}

The modulus $S$ in Persson--Volpato is the heterotic axio-dilaton, not the $T^2$
complex-structure modulus $\tau$ that a companion negative result in this
programme~\cite{Fricke2026} investigated and found \emph{not} supported, at level $N>1$, by the
natural T-duality reference (Giveon--Porrati--Rabinovici). That the same Fricke matrix
$W_N=\left(\begin{smallmatrix}0&-1\\N&0\end{smallmatrix}\right)$ recurs attached to a
\emph{different} modulus in a correctly-sourced reference, while failing to attach to the
$T^2$ modulus in an incorrectly-sourced one, is exactly the distinction that negative result was
designed to surface, and this note's citation discipline is a direct continuation of it: cite the
source that actually supports the specific claim being motivated, not the source that merely
mentions a similar-looking matrix.

\section{Related work}\label{sec:related}

Martin~\cite{Martin1996} gives the classical account of Fricke and Atkin--Lehner eigenvalue
multiplicativity for eta-quotients; Theorem~\ref{thm:main} is the special case of that classical
picture reachable directly from the already-proved general transformation law, stated and checked
in Lean rather than re-derived. \lean{etaQuotient_fricke}
(\S\ref{sec:introduction})\cite{EtaQuotients2026} supplies the general law this note specializes;
\cite{Fricke2026} supplies the matrix-level Fricke involution $W_N$ that both eta-quotient papers'
transformation laws are ultimately stated relative to, and its own T-duality negative result is the
direct precedent for this note's citation discipline in \S\ref{sec:motivation}.

\section{Limitations}\label{sec:limitations}

\begin{remark}[What is not done here]
\leavevmode
\begin{enumerate}[label=(\alph*)]
  \item No connection to $\Gamma_0(N)$-modularity, cusp behaviour, or the bundled type
    $M_k(\Gamma_0(N))$ is made (Remark~\ref{rem:not-modularity}).
  \item The physics claim motivating this note (\S\ref{sec:motivation}) is not formalized, is not
    attempted here, and should not be inferred from the Lean content.
  \item $\lambda(0,k)=0$ is junk (the definition is total but the hypothesis $N>0$ is required by
    every theorem that uses it); no theorem here relies on or claims anything about $N=0$.
  \item The characterization of $\lambda=\pm1$ (Theorem~\ref{thm:pm1}) is complete for the
    eigenvalue itself, but Corollary~\ref{cor:pm1-weight} inherits the self-duality
    \emph{hypothesis} on $r$ --- it does not characterize which exponent vectors are self-dual, a
    separate combinatorial question this note does not address beyond
    Definition~\ref{def:selfdual}'s decidability.
\end{enumerate}
\end{remark}

\section{Conclusion}

We have formalized, in Lean~4, that a Fricke-self-dual eta-quotient is an eigenform of the Fricke
transform, with a closed-form eigenvalue depending only on level and weight, and exactly when that
eigenvalue is $\pm1$. The mathematics is a short, honest corollary of an already-proved
transformation law; its interest is less in the difficulty of the proof than in the fact that the
same combinatorial condition organizes a physical duality in the CHL string-compactification
literature, reported here strictly as motivation, with the two literatures' differing
normalizations stated explicitly so they cannot be silently conflated.

\section*{A note on method: AI assistance and human responsibility}

This note follows the disclosure norm argued for by Tao~\cite{Tao2026AI}, endorsing recommendations
of the Leiden Declaration on Artificial Intelligence and Mathematics~\cite{Leiden2026}: disclose
tool use, affirm that credit and responsibility belong to the human author, and attribute sources
where AI tools cannot reliably do so themselves.

\textbf{What AI did.} The Lean~4 formalization, the literature search, and the drafting of this
text were produced with the assistance of large language models (Anthropic's Claude family),
working under human direction across multiple sessions. The methodology is neuro-symbolic in the
specific sense Tao describes~\cite{Tao2026AI}: natural-language reasoning proposes definitions,
proof strategies, and prose, and the Lean kernel --- not the model, and not the author's trust in
the model --- is what admits or rejects every mathematical claim. Every theorem in this paper
compiles against Mathlib with no \lean{sorry} and a build-failing axiom guard on its footprint
(\S\ref{sec:verification}); nothing here is asserted true because a model said so.

\textbf{What the author did.} Xavier Callens set the research direction --- choosing this
specialization specifically because it was reachable from already-proved machinery rather than
requiring a fresh physics-bridge attempt, after a companion run in this programme found such an
attempt correctly refused~\cite{Fricke2026} --- judged what the result means and whether it is
worth publishing, and is responsible for every claim in this document, including its errors. The
physics/mathematics boundary in \S\ref{sec:motivation} was treated as a hard constraint on the
formalization itself, not merely on the prose: no Lean declaration in this development names or
depends on any physical object, verified by direct inspection and independently re-verified by an
adversarial review pass. Tao's own test is apt: a result that cannot be explained and defended by
its author, formally verified or not, is not ready to be called complete~\cite{Tao2026AI}. The
mathematical content here --- why self-duality collapses the product, and why the two
normalizations in \S\ref{sec:motivation} diverge --- is something the author can explain without
further AI assistance; the Lean proof terms that certify it were substantially AI-drafted and are
trusted only because the kernel, not the author, checks them.

\begin{thebibliography}{99}
\bibitem{Fricke2026} X.~Callens, \textit{A Mathlib-Idiomatic Formalization of the Fricke Involution
and its Operator on Modular Forms}, Zenodo, 2026. \texttt{doi:10.5281/zenodo.22542571}.
\bibitem{EtaQuotients2026} X.~Callens, \textit{Eta-Quotients in Lean 4: Cusp Orders at Every Cusp,
Ligozat's Criterion at Small Level, and a Named Obstruction at General Level}, Zenodo, 2026.
\texttt{doi:10.5281/zenodo.22648098}.
\bibitem{Martin1996} Y.~Martin, \textit{Multiplicative eta-quotients}, Trans.\ Amer.\ Math.\ Soc.\ \textbf{348} (1996), 4825--4856.
\bibitem{PerssonVolpato2015} D.~Persson, R.~Volpato, \textit{Fricke S-duality in CHL models}, JHEP \textbf{12} (2015) 156. \texttt{arXiv:1504.07260}.
\bibitem{Mathlib2020} The mathlib Community, \textit{The Lean mathematical library}, Proceedings of the 9th ACM SIGPLAN International Conference on Certified Programs and Proofs (CPP), 2020.
\bibitem{deMoura2021} L.~de~Moura et al., \textit{The Lean 4 Theorem Prover and Programming Language}, CADE-28, LNCS \textbf{12699}, Springer (2021).
\bibitem{Tao2026AI} T.~Tao, \textit{Mathematics in the age of AI}, essay based on a lecture at the
International Congress of Mathematicians, July 2026, submitted to the Proceedings of the ICM 2026.
\texttt{arXiv:2608.16753}.
\bibitem{Leiden2026} \textit{The Leiden Declaration on Artificial Intelligence and Mathematics},
June 2, 2026, endorsed by the International Mathematical Union. \url{https://leidendeclaration.ai/},
\texttt{doi:10.5281/zenodo.20302944}.
\end{thebibliography}

\end{document}
